\documentclass[11pt]{jlreq}
\usepackage[a4paper,margin=1.45cm]{geometry}
\usepackage{luatexja}
\usepackage{amsmath,amssymb,mathtools,array,booktabs,longtable}
\setlength{\parindent}{1em}
\setlength{\parskip}{0pt}
\allowdisplaybreaks
\setlength{\emergencystretch}{12em}
\sloppy
\hbadness=10000
\vbadness=10000
\hfuzz=2pt
\vfuzz=10pt
\raggedbottom
\makeatletter
\renewcommand\footnotesize{\@setfontsize\footnotesize{7.5}{8.7}\selectfont}
\makeatother
\begin{document}
\providecommand{\Kfield}{\mathfrak{K}}
\providecommand{\Ssys}{\mathfrak{S}}
\providecommand{\Mfield}{\mathfrak{M}}
\providecommand{\tq}{\tau}
\providecommand{\Lsys}{\mathfrak{L}}
\providecommand{\Jdom}{\mathfrak{J}}
\providecommand{\ds}{\displaystyle}

\begin{center}
{\Large\bfseries 5. 有理関数体}\par
\vspace{1.2em}
\emph{Jahresbericht der DMV} 22 (1913), pp. 316--319
\end{center}

\vspace{1.5em}

以下の諸問題は、もともと E. Fischer との会話に由来する。なお、それらの問題のいくつかは、一つの不定元という特殊な場合について、しかも係数体に関するより一般的な仮定のもとで、すでに E. Steinitz によって提起され解決されている。\footnote{E. Steinitz: \emph{Algebraische Theorie der Körper}, とくに \S\ 24. Crelle's Journal 137, 1910.}

「有理関数体」とは、与えられた数体からの係数をもつ $n$ 個の不定元の有理関数を元とする体をいう。この係数体は、とくにすべての複素数を含んでもよい。このような体の例としては、$n$ 個の量の対称関数の体、またより一般に、ある群の置換を許す $n$ 個の不定元の有理関数全体（Lagrange の \emph{Gattungsbereiche}）、さらに不変式体がある。

最初に挙げた二つの体と最後の体との間には、特徴的な相違がある。前者は $n$ 個の代数的に独立な関数、すなわち基本対称関数を含むのに対し、代数的に独立な不変式の個数は常に不定元の個数より小さい。したがって、一般に、第二種のこのような体を、いくつかの不定元を数で置き換えることによって第一種の体と一対一に対応させられることが重要である。そこで以下では、第一種の体、すなわち $n$ 個の代数的に独立な関数をもつ体に限って考える。この対応により、結果は一般にも成り立つ。

問題は三つの基底概念、すなわち有理基底、最小基底、整性基底のまわりに集まる。

1. 「有理基底」とは、その体の有限個の関数であって、体のすべての関数が、与えられた係数体からの係数を用いて、それら有限個の関数の有理的結合として表されるものをいう。簡単な考察によって、任意の有理関数体について有理基底の存在が示される。一つの不定元の場合には、この考察は E. Steinitz の前掲書にも見られる。例えば Lagrange の定理によれば、基本対称関数と群に属する一つの関数は、先に述べた Lagrange の \emph{Gattungsbereiche} の有理基底をなす。

2. 任意の有理基底は、第一種の体に対して、少なくとも $n$ 個の関数を含まなければならない。ちょうど $n$ 個の関数を含むならば、それらは代数的に独立でなければならず、そのような基底を「最小基底」と呼ぶ。最小基底の存在問題は、Lüroth, Castelnuovo, Enriques の定理によって部分的に答えられる。\footnote{J. Lüroth: Beweis eines Satzes über rationale Kurven. \emph{Math. Ann.} 9, 1875. --- G. Castelnuovo: Sulla razionalità delle involuzioni piane. \emph{Math. Ann.} 44, 1893. --- F. Enriques: Sopra una involuzione non razionale dello spazio. \emph{Rend. Acc. Linc.} vol. 21, 1912年1月21日.} それらによれば、一つの不定元の体には常に最小基底が存在する。すなわち、その体の Lüroth 関数であり、これは一次分数変換を除いて一意に定まる（Steinitz \S\ 24 参照）。二つの不定元の体については、係数体の代数拡大を許すならば、そして一般にはその場合に限って、常に最小基底が存在する。これに対して三つ以上の不定元については、一般には最小基底はまったく存在しない。最小基底をもつ特殊な体を特徴づける試みは、まだ行われていない。

私は Lagrange の \emph{Gattungsbereiche} における最小基底の問題をさらに追究した。この場合、それは次の意味をもつ。「群 $G$ に属する Lagrange の \emph{Gattungsbereich} が最小基底をもつならば、指定された群 $G$ をもつ方程式全体を、パラメータ表示によって有理的に構成できる。しかも、最小基底の係数体を含む任意の数体についてそうであり、特にこの係数体が有理数体であるならば、任意の数体についてそうである。」

最も簡単な例は対称群である。この場合、基本対称関数は有理数係数をもつ最小基底をなす。そこからパラメータ表示として、単に不定係数をもつ方程式が得られる。これから任意の数体上の任意に多くの affektlos な方程式へ移る方法は、Hilbert の既約性定理が示している。

さて、Lüroth と Castelnuovo の定理から、3 次および 4 次方程式に現れるすべての群について最小基底の存在が直接に従う。同時に、この最小基底は有理数係数をもつことも分かる。したがって、任意の数体について、指定された群をもつ 3 次および 4 次方程式全体を有理的に構成できる。巡回群および二面体群については、$n$ 乗根を係数体とする最小基底が存在する。さらにいくつかのメタ巡回群についても同様である。最小基底がそもそもある種の群を特徴づけるかどうかは、未決のままにしておかなければならない。

3. 最後の基底概念に到達するため、体の中に含まれる多項式全体を考える。「整性基底」とは、これらの多項式の有限個であって、体の任意の多項式が、与えられた係数体からの係数を用いて、それら有限個の多項式の整有理結合として表されるものをいう。整性基底の存在問題は、例えば不変式系の有限性を特殊な場合として含むが、Hilbert の「数学の問題」\footnote{D. Hilbert: Mathematische Probleme. パリ講演, 1900. 問題14. \emph{Göttinger Nachr.} 1900.}では、やや異なる形、すなわち相対的に整な関数の問題として提出されている。

当面、私は整性基底が存在する体の一つの類だけを示すことができる。すなわち、その体が $n$ 個の多項式の系を含み、それらの最高次元項の同次終結式が零でないならば、整性基底が存在する。それは、線形形式
\[
  u_0=u_1x_1+u_2x_2+\cdots+u_nx_n\footnotemark
\]
に対する、その体で既約な方程式に現れるすべての $u$ の係数によって与えられる。
\footnotetext{この既約方程式は、一つの不定元の場合には、Steinitz による Lüroth の定理の証明に見られる。}
とくに終結式の値が係数体の単元に等しいならば、表示の整性も保証される。この類の体の例は Lagrange の \emph{Gattungsbereiche} である。なぜなら基本対称関数の終結式は値 $\pm 1$ をもつからである。さらに別の例（整性なし）は、ただ一つの代数的に独立な多項式だけを含むすべての体によって与えられる。この場合、整性基底は、すべての多項式を含む最小の部分体の Lüroth 関数と同一である。したがって、よく知られているように、$n$ 変数の二次形式のすべての整有理な射影不変式は、判別式の整関数である。

いま述べた条件は整性基底の存在に十分であるにすぎず、決して必要ではない。任意の $n$ 個の多項式の終結式が消えるにもかかわらず、あの既約方程式の係数によって与えられる整性基底をもつ体は、容易に構成できる。最も一般の場合にも、あの方程式の係数が整性基底を与えるのではないか、という問題が生じる。

\clearpage

\begin{center}
{\Large\bfseries 6. 有理関数の体と系}\par
\vspace{1.2em}
\emph{Mathematische Annalen} 76 (1915), pp. 161--196
\end{center}

\vspace{1.5em}

本論文は、有理関数および整有理関数の任意の系に関する基底問題を扱う。ここで用いられる体論の方法により、有理関数からなる体、すなわち有理関数体\footnote{ウィーン自然科学者会議の機会における予備的報告「Rationale Funktionenkörper」, \emph{Jahresber. d. D. Math.-Ver.} 22 (1913) を参照。そこでは、問題設定と関数体に関する結果の概観が与えられている。}についてこれらの問題を解くことが本質的な部分として現れ、任意の系への結果の一般化はその帰結として得られる。

一般の系に対する基底問題のうち、これまで知られていたのは Hilbert の定理（\emph{Math. Ann.} 36）によって保証される加群基底の存在だけであった。以下では、任意の系に対する有理基底の存在（\S\ 7）によって、有理的表示可能性の問題が完全に解決される。体の有理基底はすでに \S\ 4 で得られる。この有理基底の存在により、常に抽象的に定義された体または系から出発することができ、したがって、特殊な有理基底の選択だけに由来し、系そのものには由来しない困難、例えば特殊な分母や基底関数の基本点の出現を避けることができる。

体については、最小基底、すなわち代数的に独立な関数からなる有理基底の問題も扱われる（\S\ 6）。さらに体論の方法は、特別な有理基底、すなわちインボリューション基底\footnote{この名称は、幾何学的解釈において、有理関数体が線形空間内の最も一般的なインボリューションを表すことによる。}（\S\ 5）へ導く。これは多項式からなる整性領域にとって特に重要となる（\S\ 8）。その場合、固定された分母（整性領域によって定まる関数の冪）をもつ表示を与える。これは、例えば不変式の典型表示という特殊な場合に知られているものである。

有理的表示可能性から、可能な限り、整有理的表示可能性、すなわち狭い意味での有限性の問題への帰結を引き出す。

これまで知られていた有限性定理はすべて、Hilbert の加群基底定理が、表示
\[
  F=A_1f_1+A_2f_2+\cdots+A_kf_k
\]
から第二の表示
\[
  F=B_1f_1+B_2f_2+\cdots+B_kf_k
\]
が従い、ここで $B_i$ がその系に属し、かつ $F$ より低い次数をもつようなすべての系に対して有限性を保証する、という事実に基づいている。

これに対して、有理基底に関する定理と体論の方法は、まったく異なる種類の仮定のもとで有限性を推論することを可能にする。

こうして Hilbert の問題（数学の問題 14）への部分的な答えとして、相対的に整な関数の一類（第一種の相対的に整な領域）が得られる。これらは有限な整性領域であり、その整性基底は上に述べたインボリューション基底によって与えられる（\S\ 10）。さらに、Hilbert による Kronecker の代数的整数の基本系に関する定理の証明（\emph{Complete invariant systems}, \S\ 2; \emph{Math. Ann.} 42）と類似に、すべての正則多項式系、すなわち基本点をもたないものへ写すことのできるすべての系は整性基底をもつことが示される（\S\ 12）。一方、整性基底をもたない非正則系は容易に指定できる。この事実の簡単な例として、次の定理を挙げられる。「一つの不定元に関する任意の無限個の多項式の系において、その系の任意の多項式が有限個の多項式の整有理結合として表されるような有限個の多項式が存在する。」さらなる例は \S\ 13 にある。

最後に、末尾の二つの節（\S\S\ 14, 15）は、整性に関して基底定理を精密化する。

この研究の本質的な道具は \S\ 3 の転移原理である。それによれば、ここで提起されるすべての基底問題について、不定元の個数が系の代数的に独立な関数の個数と一致するような系だけを考えれば十分である。

本研究の発端は、E. Fischer 氏との会話、とくに Lagrange の \emph{Gattungsbereiche} の最小基底に関する同氏の問いであった。これに関わる研究は、指定された群をもつ方程式の構成へ導いたが（前掲の「有理関数体」第2部を参照）、それは後の公表に委ねる。

本論文が最初の報告から本質的に拡張されている点は、そこでは体または特殊な整性領域についてだけ述べられた基底定理が、ここでは任意の系へ移されていることである。この転移は、任意の系を含む最小の体という概念によって簡単に行われる。これは整性領域の商体の一般化である（\S\ 7）。私がこの概念を形成するに至ったのは、K. Hentzelt 氏が、体の有理基底を知った後、Hilbert の加群基底を経由する迂回によって任意の系の有理基底の存在を証明できた、という事実による。また、正則系の整性基底に関する定理へ私を導いたのは、整性領域に対して私が用いた推論が、同じ仮定を満たす任意の系にも有効であるという K. Hentzelt の指摘であった。さらに、いくつかの個別の小さな指摘も Hentzelt 氏に負っている。

最後に、一つの不定元の場合、体の有理基底と最小基底は E. Steinitz の \emph{Algebraische Theorie der Körper}, \S\ 24（Crelle's Journal 137）に見られることを述べておく。そこでは係数体は最も一般的な意味での体として仮定されている。これらのより一般的な仮定のもとで、ここに与える定理がどこまで保たれるかという問題は、以下では扱わない。いずれにせよ証明は修正を必要とする。例えば、関数行列式の理論からの手段は標数 $p$ の体では機能しないからである。

\section*{\S\ 1. 体 $\Kfield_{n\rho}$ と系 $\Ssys_{n\rho}$.}

以下では、$n$ 個の不定元の有理関数の系、とくに有理関数体を調べる。

$f(x_1,\ldots,x_n)$, $g(x_1,\ldots,x_n)$, \ldots、または略して $f(x)$, $g(x)$, \ldots は、常に $x_1,\ldots,x_n$ の有理関数を意味するものとする。これらは既約表示、すなわち分子と分母が互いに素である表示にあると仮定する。整有理関数（多項式）は大文字 $F(x)$, $G(x)$, \ldots で表す。係数領域 $\Omega$ は何らかの数体と仮定され、したがってとくにすべての複素数を含んでもよい。$\Omega$ はさらに有限個のパラメータを含んでもよい。

$\Omega$ からの量、したがって零次のすべての関数は、常に系に属するものと数える。

これらの取り決めのもとで、有理関数からなる体、すなわち有理関数体は抽象的に定義できる。

\emph{定義 I: 有理関数の系は、次の条件を満たすとき体と呼ばれる。}
\begin{enumerate}
\item \emph{$f(x)$ とともに、$\Omega$ からの任意の量 $c$ に対して $c\cdot f(x)$ も含む。}
\item \emph{$f(x)$ と $g(x)$ とともに、常に $f(x)+g(x)$, $f(x)\cdot g(x)$ および、$g(x)\ne0$ のとき商 $f(x):g(x)$ も含む。}\footnote{ここで $f(x)$ と $g(x)$ は零次、すなわち $\Omega$ の量であってもよい。}
\end{enumerate}
特殊な体とは、有限個の有理関数
\[
  f_1(x),\ f_2(x),\ldots, f_k(x)
\]
を $\Omega$ に添加して得られるものであり、
\[
  \Omega\bigl(f_1(x),\ldots,f_k(x)\bigr)
  \quad\text{または}\quad
  \Omega(f_1,\ldots,f_k)
\]
と表す。\footnote{これらの「特殊な体」がすでにすべての体を尽くすことは \S\ 4 で分かる。}
これら特殊な体の中には、$x_1,\ldots,x_n$ を $\Omega$ に添加して得られる体
\[
  \Omega(x_1,\ldots,x_n)
\]
も含めておく。これは $\Omega$ 係数の $x_1,\ldots,x_n$ の有理関数全体と同一である。

さらに、E. Steinitz に従って、中間体の概念を導入する。
\begin{quote}
「体 $\Mfield$ が $\Omega_1$ と $\Omega_2$ の間にあるとは、$\Mfield$ が $\Omega_1$ を含み、かつ $\Omega_2$ に含まれることをいう。」
\end{quote}
すると定義 I は次のようにも述べられる。

$n$ 個の不定元の有理関数体とは、$\Omega$ と $\Omega(x_1,\ldots,x_n)$ の間の中間体であり、少なくとも一つの関数において $n$ 個の不定元を実際に含むものである。

不定元の個数のほかに、有理関数の系にはさらに一つの特徴数、すなわち代数的に独立な関数の個数、あるいは代数的階数が、次の定義によって現れる。

\emph{定義 II: 有理関数の系が代数的階数 $\rho$ をもつとは、その系の中に代数的に独立な $\rho$ 個の関数を指定できるが、その系の任意の $(\rho+1)$ 個の関数は代数的に従属であることをいう。}\footnote{$\rho$ 個の関数 $f_1(x),\ldots,f_\rho(x)$ が代数的に独立であるとは、任意の関係
\[
F\bigl(f_1(x),\ldots,f_\rho(x)\bigr)=0\quad\text{$x_1,\ldots,x_n$ について恒等的に}
\]
から
\[
F(\lambda_1,\ldots,\lambda_\rho)=0\quad\text{$\lambda_1,\ldots,\lambda_\rho$ について恒等的に}
\]
が従うことをいう。反対の場合、それらは代数的に従属であると呼ばれる。}

$n$ 個の不定元の有理関数の（有限または無限の）系で、代数的階数 $\rho$ をもつものを、二重添字によって
\[
  \Ssys_{n\rho}
\]
と表す。特に
\[
  \Kfield_{n\rho}
\]
は、$n$ 個の不定元と代数的階数 $\rho$ をもつ体を意味する。次の不等式が成り立つ。
\[
  1\leq \rho\leq n.
\]

体 $\Kfield_{nn}$ および $\Kfield_{n\rho}$ の例をさらにいくつか挙げる。

1. 体 $\Kfield_{nn}$ として、Lagrange の \emph{Gattungsbereiche} がある。これは
\begin{quote}
「$x_1,\ldots,x_n$ の置換群 $G$ の置換を許す、$x_1,\ldots,x_n$ の有理関数（および整有理関数）全体」
\end{quote}
として定義できる。これらは常に対称関数の体を含み、したがって $n$ 個の代数的に独立な基本対称関数も含む。

2. 体 $\Kfield_{n\rho}$ として、射影不変式体がある。これは一つまたは複数の基本形式の有理不変式（および整有理不変式）全体から作られる。\footnote{行列式 $1$ の変換だけを考えるのが便利である。そのとき任意個の不変式の任意の有理結合も再び変換を許し、実際に体が得られる。斉次・等重の不変式だけを取り出したものは体をなさないが、系 $\Ssys_{n\rho}$ をなす。} ここでは $\rho<n$ である。不変式の代数的に独立な個数は、常に係数の個数より小さいからである。

射影群の部分群に関する不変式体についても、対応する主張が成り立つ。

\section*{\S\ 2. 代数的に従属な関数に関する補助定理.}

この節の補助定理により、\S\ 3 で、系の代数的階数がその本来的な特徴数であることが示される。すなわち補助定理は、$\rho$ 個の代数的に独立な関数が代数的に独立なままでいるように、いくつかの不定元を数値によって逐次特殊化しても、それら $\rho$ 個の関数に代数的に従属な任意の関数は、恒等的に消えることも不定になることもない、ということを示す。

したがって、系
\begin{equation}
  g_1(x),\ g_2(x),\ldots,g_\rho(x)
\tag{1}
\end{equation}
が代数的階数 $\rho$ をもち、ここで $\rho<n$ とし、$h(x)$ を関数 (1) に代数的に従属な有理関数とする。既知の定理により、(1) の関数行列の階数は $\rho$ に等しい。したがって不定元は、例えば次のように名づけられる。
\begin{equation}
 \frac{\partial(g_1,g_2,\ldots,g_\rho)}{\partial(x_1,x_2,\ldots,x_\rho)}
 =\frac{\Gamma(x)}{G(x)^{\rho+1}}\ne0,
\tag{2}
\end{equation}
ここで $G(x)$ は通常どおり、関数 (1) の共通分母を意味する。

\emph{いま数 $a$ を、積 $G(x)\cdot\Gamma(x)$ が $(x_i-a)$ で割り切れないように選ぶ}。ただし $i$ は $\rho+1,\rho+2,\ldots,n$ のいずれかとする。したがって $x_i=a$ としても、関数 (1) の共通分母は消えず、$\rho$ 個の関数
\begin{equation}
  [g_1(x)]_{x_i=a},\ldots,[g_\rho(x)]_{x_i=a}
\tag{3}
\end{equation}
もまた代数的に独立である。言い換えれば、系 (1) の代数的階数は、置換 $x_i=a$ のもとで保存される。

関数 (1) に関して $h(x)$ が満たす既約方程式を
\begin{equation}
\begin{aligned}
&A_0(g_1,\ldots,g_\rho)\,h(x)^\tau
 +A_1(g_1,\ldots,g_\rho)\,h(x)^{\tau-1}
 +\cdots \\
&\hspace{6em}+A_\tau(g_1,\ldots,g_\rho)=0
\end{aligned}
\tag{4}
\end{equation}
とする。ただし $A_0(g_1,\ldots,g_\rho)$ および $A_\tau(g_1,\ldots,g_\rho)$ は恒等的には消えない。これを多項式の間の恒等式に移すために、
\begin{equation}
  g_i(x)=G_i(x):G(x);\qquad h(x)=H_1(x):H_2(x)
\tag{5}
\end{equation}
と置く。そして、$h(x)$ が既約表示にあるという仮定により、$H_1(x)$ と $H_2(x)$ は互いに素であることに注意する。(5) によって (4) は
\begin{equation}
\begin{aligned}
&B_0(G,G_1,\ldots,G_\rho)\,G(x)^{\sigma_0}H_1(x)^\tau \\
&\quad+B_1(G,G_1,\ldots,G_\rho)\,G(x)^{\sigma_1}H_1(x)^{\tau-1}H_2(x)+\cdots \\
&\quad+B_\tau(G,G_1,\ldots,G_\rho)\,G(x)^{\sigma_\tau}H_2(x)^\tau=0
\end{aligned}
\tag{6}
\end{equation}
となる。ここで $B_k$ は $G,G_i$ の同次形式であり、非負指数 $\sigma_k$ の少なくとも一つは零でなければならない。

いま $H_1(x)$ が $(x_i-a)$ で割り切れると仮定する。すると (6) により
\[
  B_\tau(G,G_1,\ldots,G_\rho)\,G(x)^{\sigma_\tau}H_2(x)^\tau
\]
も $(x_i-a)$ で割り切れる。これは $G(x)$ についての仮定により排除されている。したがって $B_\tau$ の割り切れ性から
\[
  A_\tau=B_\tau:G^\lambda
\]
の割り切れ性も従い、関数 (3) の間に関係 $A_\tau=0$ が存在することになり、仮定された代数的独立性に反する。最後に、$H_2(x)$ は $H_1(x)$ と互いに素であるから、これも $(x_i-a)$ で割り切れることはできない。\footnote{既約表示に基づくこの最後の結論は、残りの仮定が保たれるとしても、置換 $x_i=a$, $x_k=b$ に対してはもはや可能ではない。その場合 $H_1(x)$ と $H_2(x)$ は同時に消えることがあり、商はその置換のもとで $\frac{0}{0}$ に等しい不定形となりうる。例えば
\[
 h(x)=\frac{x_3(\alpha x_1+\beta x_2)+x_4(\gamma x_1+\delta x_2)}
 {x_3(\alpha' x_1+\beta' x_2)+x_4(\gamma' x_1+\delta' x_2)}
\]
を $x_3=0$, $x_4=0$ と置換する場合である。} よって $H_1(x)$ が $(x_i-a)$ で割り切れるという仮定は矛盾に至る。同様に、$H_2(x)$ も $(x_i-a)$ で割り切れないことが示される。したがって関数 $[h(x)]_{x_i=a}=h'(x)$ は、恒等的に消えることも不定になることもない。

以上の結論を、再び既約表示に置いた系 (3) と関数 $h'(x)$ に適用できる。この手続きを続けることで、最終的に次の補助定理に到達する。

\emph{補助定理: 二つの系}
\[
  g_1(x),\ g_2(x),\ldots,g_\rho(x),
\]
\[
  g_1(x),\ g_2(x),\ldots,g_\rho(x),\ h(x)
\]
\emph{はいずれも代数的階数 $\rho$ をもつものとし、ここで $\rho<n$ である。例えば}
\[
 \frac{\partial(g_1,\ldots,g_\rho)}{\partial(x_1,\ldots,x_\rho)}
 =\frac{\Gamma(x)}{G(x)^{\rho+1}}\ne0
\]
\emph{と仮定する。数}
\[
  a_n,\ a_{n-1},\ldots,a_{\rho+1}
\]
\emph{は、置換}
\[
  x_n=a_n,\quad x_{n-1}=a_{n-1},\ldots,\quad x_{\rho+1}=a_{\rho+1}
\]
\emph{のもとで積 $G(x)\cdot\Gamma(x)$ が消えないように、すなわち系 (1) の代数的階数が保存されるように選ばれるものとする。逐次置換}
\[
\begin{aligned}
 [h(x)]_{x_n=a_n}&=h^{(n-1)}(x),\qquad
 [h^{(n-1)}(x)]_{x_{n-1}=a_{n-1}}=h^{(n-2)}(x),\ldots,\\
 [h^{(\rho+1)}(x)]_{x_{\rho+1}=a_{\rho+1}}&=h^*(x_1,x_2,\ldots,x_\rho)
\end{aligned}
\]
\emph{において、関数 $h(x)$, $h^{(n-1)}(x)$, \ldots, $h^{(\rho+1)}(x)$ をそれぞれ既約表示に置く。これらの仮定のもとで、関数 $h^*(x_1,\ldots,x_\rho)$ では分子も分母も恒等的には消えず、また関数は $\frac{0}{0}$ になることもない。}\footnote{例えば前注の関数 $h(x)$ では、逐次置換により
\[
 [h(x)]_{x_4=0}=h^{(3)}(x)=\frac{\alpha x_1+\beta x_2}{\alpha' x_1+\beta' x_2},\qquad
 h^*(x_1,x_2)=\frac{\alpha x_1+\beta x_2}{\alpha' x_1+\beta' x_2}
\]
が得られる。}

\section*{\S\ 3. 系 $\Ssys_{n\rho}$ から系 $\Ssys^*_{\rho\rho}$ への一対一対応.}

補助定理から、いくつかの不定元を数で置き換えることにより、系 $\Ssys_{n\rho}$ を系 $\Ssys^*_{\rho\rho}$ へ対応させることが直ちに得られる。このとき代数的階数 $\rho$ は本来的な特徴数として現れる。

$\Ssys_{n\rho}$ は代数的階数 $\rho$ をもつので、$\Ssys_{n\rho}$ から $\rho$ 個の関数
\begin{equation}
  g_1(x),\ldots,g_\rho(x)
\tag{1}
\end{equation}
が存在し、それらは代数的に独立である。さらに、$f(x)$ が $\Ssys_{n\rho}$ からの任意の関数を表すならば、系
\begin{equation}
  g_1(x),\ldots,g_\rho(x),\ f(x)
\tag{2}
\end{equation}
もまた代数的階数 $\rho$ をもち、系 (1) と (2) は補助定理の条件を満たす。\footnote{必要ならば、不定元を番号づけ直す。}

したがって、数
\[
  a_n,\ a_{n-1},\ldots,a_{\rho+1}
\]
を補助定理に従って、系 (1) の代数的階数 $\rho$ が置換によって保存されるように定める。このことは任意に多くの仕方で可能である。そのとき
\[
\begin{aligned}
 [f(x)]_{x_n=a_n}&=f^{(n-1)}(x),\qquad
 [f^{(n-1)}(x)]_{x_{n-1}=a_{n-1}}=f^{(n-2)}(x);\ldots,\\
 [f^{(\rho+1)}(x)]_{x_{\rho+1}=a_{\rho+1}}&=f^*(x_1,\ldots,x_\rho)
\end{aligned}
\]
で定義される関数、すなわち
\begin{equation}
  f^*(x_1,\ldots,x_\rho)
\tag{3}
\end{equation}
では、分子も分母も恒等的に消えることができない。

いま $f(x)$ に $\Ssys_{n\rho}$ のすべての（有限または無限個の）関数を走らせる。すべての関数 (3) 全体からなる系を考えると、それは次の性質をもつ。

1. 代数的階数 $\rho$ をもつ。実際、それは (1) から生じる関数
\[
  g_1^*(x_1,\ldots,x_\rho),\ldots,g_\rho^*(x_1,\ldots,x_\rho)
\]
を含み、これらは $a_n,a_{n-1},\ldots,a_{\rho+1}$ の選び方により代数的に独立である。したがってこの系を
\[
  \Ssys^*_{\rho\rho}
\]
と表す。\footnote{対応して、$f^*(x)$, $g^*(x)$, \ldots は、以下いつでも対応関数、すなわち $x_1,\ldots,x_\rho$ の関数を意味するものとする。}

2. 関数 $f(x)$ と $f^*(x)$ の対応は例外なく一対一である。

実際、$f_1(x)$ と $f_2(x)$ が $\Ssys_{n\rho}$ に属すならば、その差
\[
  d(x)=f_1(x)-f_2(x)
\]
も系 (1) に代数的に従属であり、さらに
\[
  d^*(x)=f_1^*(x)-f_2^*(x)
\]
である。$d(x)$ は系 (1) とともに補助定理の条件を満たすので、
\[
  d(x)\ne0
\]
から必ず
\[
  d^*(x)\ne0
\]
が従う。すなわち、$\Ssys_{n\rho}$ の異なる関数には、$\Ssys^*_{\rho\rho}$ の異なる関数が対応する。

3. $\Ssys^*_{\rho\rho}$ の関数の間の任意の関係
\[
  F\bigl(f_1^*(x),f_2^*(x),\ldots,f_k^*(x)\bigr)=0
  \quad\text{$x_1,\ldots,x_\rho$ について恒等的に}
\]
から、一意に対応する $\Ssys_{n\rho}$ の関数について
\[
  F\bigl(f_1(x),f_2(x),\ldots,f_k(x)\bigr)=0
  \quad\text{$x_1,\ldots,x_n$ について恒等的に}
\]
が従う。なぜなら、$f_1(x),f_2(x),\ldots,f_k(x)$ とともに、これらの任意の有理結合も系 (1) に代数的に従属だからである。

さらにこれは
\begin{equation}
\begin{aligned}
 h(x)&=F\bigl(f_1(x),\ldots,f_k(x)\bigr):\\
 h^*(x)&=F\bigl(f_1^*(x),\ldots,f_k^*(x)\bigr)
\end{aligned}
\tag{4}
\end{equation}
からも従う。補助定理により、したがって
\[
  h(x)\ne0\quad\text{ならば}\quad h^*(x)\ne0
\]
である。証明終わり。

要約すれば次を得る。

\emph{定理 I: 有理関数の（有限または無限の）任意の系 $\Ssys_{n\rho}$ に対して、任意に多くの仕方で、一対一の系 $\Ssys^*_{\rho\rho}$ を割り当てることができる。その際、$\Ssys^*_{\rho\rho}$ の関数間に存在するすべての関係は $\Ssys_{n\rho}$ においても保存され、また逆も成り立つ。特に (4) により、体 $\Kfield_{n\rho}$ には再び体 $\Kfield^*_{\rho\rho}$ が対応する。}

定理 I から、以後のすべての証明を簡単にする帰結を引き出す。\emph{系の関数間の有限または無限個の関係によって特徴づけられる系 $\Ssys_{n\rho}$ の性質は、対応する系 $\Ssys^*_{\rho\rho}$ について成り立てば、もとの系 $\Ssys_{n\rho}$ についても成り立つ。}

この帰結を略して「転移原理」と呼ぶ。

\providecommand{\Lsys}{\mathfrak{L}}
\providecommand{\Jdom}{\mathfrak{J}}

\clearpage
この対応の最も簡単な例は、代数的不変式の理論によって与えられる。実際、線形変換によって、基本形式のいくつかの係数を数値的に固定し、特殊な形式に対して成り立つすべての関係が一般の場合にも成り立つようにすることが常にできる。

\section*{\S\ 4. 体 $\Kfield_{n\rho}$ の有理基底.}

以下の研究では、基底概念のまわりにまとめられた問題に対して、\S\ 3 の「転移原理」を一貫して用いる。まず体について基底の存在を証明するが、定義は一般に述べる。

\emph{定義 III: $\Ssys_{n\rho}$ の有限個の関数が有理基底と呼ばれるのは、$\Ssys_{n\rho}$ の任意の関数が、$\Omega$ からの係数を用いて、それら有限個の関数の有理結合として表されるときである。}

有理基底は $\rho$ 個の代数的に独立な関数を含まなければならない。なぜなら、有理基底から導かれる系の代数的階数は、有理基底の代数的階数に等しいからである。

まず、有理基底の存在証明が転移原理を許すことを示す。すなわち、定義 III は次のように定式化できる。

$\Ssys_{n\rho}$ の関数
\[
  g_1(x),\ g_2(x),\ldots, g_k(x)
\]
が有理基底をなすことと、無限個の関係
\begin{equation}
  f(x)=\gamma\bigl(g_1(x),\ldots,g_k(x)\bigr)
\tag{1}
\end{equation}
が満たされることは同値である。ここで $f(x)$ は $\Ssys_{n\rho}$ のすべての関数を走り、$\gamma$ は $f$ に応じて変わる $k$ 個の引数の有理関数を表す。

したがって、対応系 $\Ssys^*_{\rho\rho}$ において
\[
  g_1^*(x),\ g_2^*(x),\ldots,g_k^*(x)
\]
が有理基底を与え、それにより関係
\[
  f^*(x)=\gamma\bigl(g_1^*(x),\ldots,g_k^*(x)\bigr)
\]
が生じるならば、定理 I によって、関係 (1) は $\Ssys_{n\rho}$ に対して成り立つ。したがって、\emph{有理基底に対する転移原理}は次のようにいう。

\emph{対応系 $\Ssys^*_{\rho\rho}$ に有理基底が存在することから、もとの系 $\Ssys_{n\rho}$ の有理基底が従う。}\footnote{転移原理は、ここで扱うすべての基底問題について同様に定式化される。}

すべての体 $\Kfield_{n\rho}$ の有理基底の存在を証明するためには、したがって体 $\Kfield^*_{\rho\rho}$ に限ればよい。ここでは E. Steinitz が用いた議論の直接の一般化が目的に至る。\footnote{E. Steinitz, \emph{Algebraische Theorie der Körper}, \S\ 24, Crelle's Journal 137 (1909).}

\begin{equation}
  g_1^*(x_1,\ldots,x_\rho),\ldots,
  g_\rho^*(x_1,\ldots,x_\rho)
\tag{2}
\end{equation}
を、$\Kfield^*_{\rho\rho}$ からの $\rho$ 個の代数的に独立な関数の系とする。$x_1,\ldots,x_\rho$ を消去すると、$\rho$ 個の関係
\begin{equation}
\begin{aligned}
  F_1\bigl(g_1^*(x),\ldots,g_\rho^*(x),x_1\bigr)&=0,\ldots,\\
  F_\rho\bigl(g_1^*(x),\ldots,g_\rho^*(x),x_\rho\bigr)&=0
\end{aligned}
\tag{3}
\end{equation}
が得られる。各関係 $F_i=0$ には不定元 $x_i$ が実際に現れる。そうでなければ、仮定に反して関数 (2) の間に代数的従属性が存在することになる。関係 (3) は
\[
  x_1,x_2,\ldots,x_\rho
\]
が
\[
  \Omega\bigl(g_1^*(x),\ldots,g_\rho^*(x)\bigr)
\]
に関して代数的元であることをいう。したがって、$x_1,\ldots,x_\rho$ を添加して得られる体
\[
  \Omega\bigl(g_1^*,\ldots,g_\rho^*,x_1,
  \ldots,x_\rho\bigr)=\Omega(x_1,
  \ldots,x_\rho)
\]
は、$\Omega(g_1^*,\ldots,g_\rho^*)$ の有限代数拡大である。既知の定理\footnote{例えば Weber, \emph{Lehrbuch d. Algebra} 1, \S\ 144（第1版）または \S\ 55（小版）を比較せよ。}により、$\Kfield^*_{\rho\rho}$ は、$\Omega(g_1^*,\ldots,g_\rho^*)$ と $\Omega(x_1,\ldots,x_\rho)$ の間の中間体として、それ自身 $\Omega(g_1^*,\ldots,g_\rho^*)$ の代数拡大であり、また $\Omega(x_1,\ldots,x_\rho)$ は $\Kfield^*_{\rho\rho}$ の代数拡大である。したがって $\Kfield^*_{\rho\rho}$ は、$\Omega(g_1^*,\ldots,g_\rho^*)$ に $\Kfield^*_{\rho\rho}$ からの有限個の元を添加することによって、または適当に選ばれた一つの関数を添加することによって得られる。すなわち有理基底が存在する。転移原理によって次を得る。

\emph{定理 II: 任意の体 $\Kfield_{n\rho}$ は有理基底をもつ。}

\section*{\S\ 5. 体 $\Kfield_{n\rho}$ のインボリューション形式とインボリューション基底.}

前の議論から、体そのものによって定まる有理基底をさらに得る。まず体 $\Kfield^*_{\rho\rho}$ についてそれを示す。

$\Kfield^*_{\rho\rho}$ を $\Omega$ と $\Omega(x_1,\ldots,x_\rho)$ の間の体とする。\S\ 4 と同様、$\Kfield^*_{\rho\rho}$ の中から $\rho$ 個の代数的に独立な関数
\[
  g_1^*(x),\ldots,g_\rho^*(x)
\]
を選び、$\Omega(g_1^*,\ldots,g_\rho^*)$ の有限代数拡大 $\Omega(x_1,\ldots,x_\rho)$ を作る。$\Omega(x_1,\ldots,x_\rho)$ の $\Kfield^*_{\rho\rho}$ に関する次数を $i$ とする。

不定係数 $u_1,\ldots,u_\rho$ をもつ線形形式
\[
  u(x)=u_1x_1+\cdots+u_\rho x_\rho
\]
を導入し、$u(x)$ の $\Kfield^*_{\rho\rho}$ に関する既約方程式を
\begin{equation}
 \Phi^*(z,u)=z^i+\sum_{\alpha+\alpha_1+\cdots+\alpha_\rho=i}
 g^*_{\alpha\alpha_1\ldots\alpha_\rho}(x)
 z^\alpha u_1^{\alpha_1}\cdots u_\rho^{\alpha_\rho}
\tag{1}
\end{equation}
とする。\footnote{Noether は総和を総和 $i$ をもつすべての添字系について書いている。ここでは主項 $z^i$ を分けている。} 係数
\[
  g^*_{\alpha\alpha_1\ldots\alpha_\rho}(x)
\]
は $\Kfield^*_{\rho\rho}$ に属し、体によって定まる。$\Phi^*(z,u)$ を体のインボリューション形式と呼ぶ。

これらの係数が有理基底をなすことを主張する。実際、(1) を $z$ の方程式として見れば、それは $\Kfield^*_{\rho\rho}$ に関して既約であり、したがって部分体
\[
  \Omega\bigl(g^*_{\alpha\alpha_1\ldots\alpha_\rho}(x)\bigr)
\]
に関してはなおさら既約である。ゆえに $\Phi^*(z,u)$ は、この部分体に関する $u(x)$ の既約方程式を与える。したがって $\Omega(x_1,\ldots,x_\rho)$ は、$\Omega(g^*_{\alpha\alpha_1\ldots\alpha_\rho}(x))$ の代数拡大、しかも次数 $i$ の拡大である。$\Kfield^*_{\rho\rho}$ はこの部分体と $\Omega(x_1,\ldots,x_\rho)$ の間にあるから、次数の一致は体の一致を意味することが知られている。転移原理により、$\Phi(z,u)$ の係数は対応して $\Kfield_{n\rho}$ の有理基底を与える。すなわち次を得る。

\emph{定理 III: インボリューション形式の係数は、体によって定まる有理基底、すなわちインボリューション基底をなす。体 $\Kfield^*_{\rho\rho}$ については、この基底は一意に定まる。}\footnote{体 $\Kfield_{n\rho}$ については、インボリューション形式は対応の選択に依存しうる。不定係数をもつ対応を用いれば一意性を得ることができる。}

この事実に非有理的な解釈を付け加える。これはインボリューションの幾何学的理論とのつながりを与える。

$\Phi^*(z,u)$ の一次因子への分解は、拡大体において
\begin{equation}
\begin{aligned}
 \Phi^*(z,u)&=(z-u_1x_1-\cdots-u_\rho x_\rho)
 (z-u_1x_1^{(1)}-\cdots-u_\rho x_\rho^{(1)})\cdots\\
 &\qquad\cdot
 (z-u_1x_1^{(i-1)}-\cdots-u_\rho x_\rho^{(i-1)})
\end{aligned}
\tag{4}
\end{equation}
によって与えられる。(1) と比較すれば、関数
\[
  g^*_{\alpha\alpha_1\ldots\alpha_\rho}(x_1,
  \ldots,x_\rho)\quad(\alpha+\alpha_1+\cdots+\alpha_\rho=i)
\]
は、量の行
\begin{equation}
\begin{array}{cccc}
  x_1 & x_2 & \cdots & x_\rho\\
  x_1^{(1)} & x_2^{(1)} & \cdots & x_\rho^{(1)}\\
  \vdots & \vdots & & \vdots\\
  x_1^{(i-1)} & x_2^{(i-1)} & \cdots & x_\rho^{(i-1)}
\end{array}
\tag{5}
\end{equation}
の基本対称関数と同一であることが分かる。したがって $\Kfield^*_{\rho\rho}$ の任意の関数は、行 (5) に関して対称なこれらの基本関数の有理結合であり、逆に行 (5) の任意の対称関数は $\Kfield^*_{\rho\rho}$ に属する。

したがって
\[
  f^*(x)=\frac{F^*(x)}{H^*(x)}
\]
が $\Kfield^*_{\rho\rho}$ のすべての関数を走るとき、無限個の関係の系
\begin{equation}
  \frac{F^*(x)}{H^*(x)}=
  \frac{F^*(x^{(1)})}{H^*(x^{(1)})}=\cdots=
  \frac{F^*(x^{(i-1)})}{H^*(x^{(i-1)})}
\tag{6}
\end{equation}
が存在する。要約すれば、
\[
  F^*(x^{(k)})H^*(x^{(l)})-
  F^*(x^{(l)})H^*(x^{(k)})=0
  \quad(k,l=0,1,\ldots,i-1),
\]
である。ここで $F^*(x^{(k)})$ も $H^*(x^{(k)})$ も恒等的には消えない。これらは $F^*(x)$ および $H^*(x)$ の共役であり、形式的に対称だからである。

逆に、無限個の関係
\begin{equation}
  F^*(x)H^*(\xi)-H^*(x)F^*(\xi)=0,
  \qquad H^*(\xi)\ne0
\tag{7}
\end{equation}
を満たす任意の量の行 $\xi$ について、量の行 (5) への所属が得られる。

実際、$G^*(x)$ を $\Phi^*(z,u)$ の係数の共通分母とすると、仮定 (7) により関数
\[
  G^*(\xi)\cdot\Phi^*(z,u;\xi)
\]
は、$\xi$ に関しても整であり、恒等的には消えない。すなわち $\Phi^*(z,u;\xi)$ の係数は不定にならず、かつ零点
\[
  z=u_1\xi_1+u_2\xi_2+\cdots+u_\rho\xi_\rho
\]
をもつ。(7) により、したがって $\Phi^*(z,u;x)$ も零点 $z=u(\xi)$ をもつ。すなわち $\xi$ は量の行 (5) に属する。$\xi$ が共役な量の行に関する方程式系
\[
  F^*(x^{(k)})H^*(\xi)-H^*(x^{(k)})F^*(\xi)=0;
  \qquad H^*(\xi)\ne0
\]
を満たす場合にも、(6) により対応する主張が成り立つ。

したがって、量の行 (5) は、方程式系
\[
  F^*(x)H^*(\xi)-H^*(x)F^*(\xi)=0;
  \qquad H^*(\xi)\ne0
\]
の解 $\xi$ として定義される。ここで $F^*(x):H^*(x)$ は $\Kfield^*_{\rho\rho}$ のすべての関数を走る。また行 $x$ は任意の共役な行で置き換えてもよい。

幾何学的には、各行 $x$ を点と解釈すると、これは次のことを意味する。(7) の解は、$\rho$ 次元線形空間において、各々 $i$ 個の点からなる $\infty^\rho$ 個の群を表し、群の任意の点がその個々の群を一意に定める。このような点群の系を「\emph{線形空間におけるインボリューション}」と呼ぶ。\footnote{(7) の除外された解、すなわち $F^*(\xi)=0$, $H^*(\xi)=0$ となるもの、およびインボリューション基底に対応する方程式系の除外された解は、インボリューションの基本点と呼ばれる。これには同次化によって生じるもの、すなわち「無限遠にある」点も含まれる。}

対応して、体 $\Kfield_{n\rho}$ は、曲線、曲面などからなる線形空間内のインボリューションを特徴づける。

以上により、$\Omega(x_1,\ldots,x_\rho)$ の $\Kfield^*_{\rho\rho}$ に関する次数 $i$ は、側条件 $H^*(\xi)\ne0$ を満たす (7) の解系の個数に等しい。これを $\Kfield^*_{\rho\rho}$ の解系と呼ぶことができる。したがって、インボリューション基底の存在を証明するために用いた議論は、次の非有理的な定式化も許す。

「$\mathfrak{L}^*$ が $\Kfield^*_{\rho\rho}$ の約体であり、体 $\mathfrak{L}^*$ のすべての解系が $\Kfield^*_{\rho\rho}$ の解系でもあるならば、$\mathfrak{L}^*$ と $\Kfield^*_{\rho\rho}$ は同一である。」

対応する命題は、転移原理により、$\Kfield_{n\rho}$ の約体 $\mathfrak{L}$ についても成り立つ。

\section*{\S\ 6. 体 $\Kfield_{n\rho}$ の最小基底.}

この節では既知の定理を概観する。ただし、転移原理と有理基底の存在によって可能となった拡張された定式化で述べる。

\emph{定義 IV: $\Kfield_{n\rho}$ の有理基底が最小基底と呼ばれるのは、それが最小個数 $\rho$ の関数だけを含むときであり、その場合それらは代数的に独立でなければならない。}

Lüroth, Castelnuovo, Enriques の定理\footnote{J. Lüroth: Beweis eines Satzes über rationale Kurven, \emph{Math. Ann.} 9 (1875). --- G. Castelnuovo: Sulla razionalità delle involuzioni piane, \emph{Math. Ann.} 44 (1893). --- F. Enriques: Sopra una involuzione non razionale dello spazio, \emph{Rend. Acc. Linc.}, vol. 21, 1912年1月21日.}から次の事実が得られる。

I. 任意の体 $\Kfield_{n1}$ は最小基底をもつ。それはその体の Lüroth 関数であり、一次分数変換を除いて一意に定まる。

II. 体 $\Kfield_{n2}$ については、係数体 $\Omega$ の代数拡大を許すならば、そして一般にはその場合に限って、常に最小基底が存在する。

III. 体 $\Kfield_{n\rho}$（$\rho\ge3$）については、一般には最小基底は存在しない。最小基底をもつ特殊な体 $\Kfield_{n\rho}$ の特徴づけはまだ試みられていない。

E. Steinitz が与えた体 $\Kfield^*_{11}$ に対する Lüroth の定理の証明を簡単に示す。\footnote{前掲書, \S\ 24.}

$\Phi^*(z,u)$ を $\Kfield^*_{11}$ のインボリューション形式とし、$\psi^*(x)$ を、実際に不定元 $x$ を含む $\Phi^*(z,u)$ の任意の係数とする。すると、$\Omega(x)$ の $\Omega(\psi^*(x))$ に関する次数は、$\Omega(x)$ の $\Kfield^*_{11}$ に関する次数に等しいことが分かる。そこから体の同一性が従う。さらに、$\Omega(\chi^*(x))$ が $\Omega(\psi^*(x))$ に等しいならば、$\chi^*(x)$ と $\psi^*(x)$ の相互の有理従属性は、二つの関数の一次分数従属性を意味する。

したがって、転移原理により、Lüroth の定理は次の形でも述べられる。

\emph{体 $\Kfield_{n1}$ の最小基底は、インボリューション基底の任意の一つの関数からなる。また $\Kfield_{n1}$ のインボリューション基底の任意の二つの関数は、一次分数変換によって互いに依存する。}

G. Castelnuovo は、幾何学的なインボリューション理論を用いて、有理基底が三つの関数から成る体 $\Kfield^*_{22}$ について事実 II を証明している。転移原理は係数体 $\Omega$ の有限代数拡大のもとでも有効であるから、有理基底の存在によって、すべての体 $\Kfield_{n2}$ について証明が完結する。

III については、Enriques が三次元空間内の非有理的インボリューション、すなわち最小基底をもたない体 $\Kfield_{33}$ を構成していることを注意しておく。
\end{document}
